\begin{lemma}[Cone comparison and radial monotonicity]
Let $E$ be a finite-dimensional real inner-product space, let $D\subseteq E$,
let $w:E\to\R$ be infinity-harmonic on $D$ in the sense of
\noderef{IsInfHarmonicOn}, let $z\in E$ and let $r>0$ with
$\overline B_r(z)\subseteq D$. Then, with the cone slopes of
\noderef{ConeSlopeUpper} and \noderef{ConeSlopeLower},
\begin{align}
  w(Z)-w(z)&\le L_r^+(w;z)\,|Z-z| &&(Z\in\overline B_r(z)),
     \label{eq:cone-upper}\\
  -\bigl(w(Z)-w(z)\bigr)&\le L_r^-(w;z)\,|Z-z| &&(Z\in\overline B_r(z)),
     \label{eq:cone-lower}
\end{align}
and for every $r'\ge r$ with $\overline B_{r'}(z)\subseteq D$,
\begin{equation}\label{eq:cone-monotone}
  L_r^+(w;z)\le L_{r'}^+(w;z),\qquad
  L_r^-(w;z)\le L_{r'}^-(w;z).
\end{equation}
\end{lemma}
\begin{proof}
Throughout, $\overline B_\rho(z)=\{Z:|Z-z|\le\rho\}$.

\emph{Step 0: the slopes are finite and nonnegative.}
Since $w$ is infinity-harmonic on $D$, it is continuous on $D$ by
\noderef{IsInfHarmonicOn}; as $\overline B_r(z)\subseteq D$ is compact, the set
$\{r^{-1}(w(z+h)-w(z)):|h|\le r\}$ is a nonempty bounded set of reals, so
$L_r^+(w;z)$ is its maximum. Taking $h=0$ shows $L_r^+(w;z)\ge0$. The same
applies to $L_r^-(w;z)$.

\emph{Step 1: $-w$ is infinity-harmonic on $D$.}
Let $\phi$ be $C^2$ near a point $Z_0\in D$ and suppose
$Z\mapsto(-w)(Z)-\phi(Z)$ has a local maximum on $D$ at $Z_0$. Then
$Z\mapsto w(Z)-(-\phi)(Z)$ has a local minimum on $D$ at $Z_0$, and $-\phi$ is
$C^2$ near $Z_0$, so the supersolution inequality of
\noderef{IsInfHarmonicOn} gives $\Dinf(-\phi)(Z_0)\le0$. By the definition of
$\Dinf$ in \noderef{InfLap}, $\nabla(-\phi)=-\nabla\phi$ and
$D^2(-\phi)=-D^2\phi$, so
\[
  \Dinf(-\phi)(Z_0)=-D^2\phi(Z_0)[-\nabla\phi(Z_0)]
   =-D^2\phi(Z_0)[\nabla\phi(Z_0)]=-\Dinf\phi(Z_0),
\]
using that $A[-\xi]=A[\xi]$ for a bilinear form $A$. Hence
$\Dinf\phi(Z_0)\ge0$, which is the subsolution inequality for $-w$. The
symmetric computation with local minima gives the supersolution inequality for
$-w$, and $-w$ is continuous on the open set $D$. Thus $-w$ is
infinity-harmonic on $D$. Note also that, directly from
\noderef{ConeSlopeUpper} and \noderef{ConeSlopeLower},
\begin{equation}\label{eq:cone-neg}
  L_\rho^+(-w;z)=L_\rho^-(w;z),\qquad
  L_\rho^-(-w;z)=L_\rho^+(w;z)
\end{equation}
for every $\rho$, because the defining sets of quotients are exchanged.

\emph{Step 2: the perturbed cone barrier.}
We prove \eqref{eq:cone-upper}. Fix any $L>L_r^+(w;z)$ and choose $\gamma>0$ so
small that
\[
  \gamma r<L-L_r^+(w;z),\qquad 2\gamma r<L .
\]
Define, for $Z\in E$,
\[
  \Theta(Z)=w(z)+L|Z-z|-\gamma|Z-z|^2 .
\]
On the sphere $|Z-z|=r$ we have, using Step 0,
\[
  w(Z)\le w(z)+rL_r^+(w;z)<w(z)+r\bigl(L-\gamma r\bigr)=\Theta(Z),
\]
the strict inequality coming from $\gamma r<L-L_r^+(w;z)$. Thus
$w-\Theta<0$ on $\{|Z-z|=r\}$, while $(w-\Theta)(z)=0$.

Suppose, for contradiction, that $w(Z)>\Theta(Z)$ for some
$Z\in\overline B_r(z)$. The function $w-\Theta$ is continuous on the compact set
$\overline B_r(z)$, so it attains a maximum there, and that maximum is strictly
positive. By the previous paragraph the maximum is not attained on the sphere,
and it is not attained at $z$ since the value there is $0$. Let $Z_0$ be a
maximum point; then $|Z_0-z|=:\rho\in(0,r)$ and $Z_0$ is interior to
$\overline B_r(z)$.

The function $\Theta$ is $C^2$ on $E\setminus\{z\}$: indeed
$Z\mapsto|Z-z|^2$ is a polynomial, and $Z\mapsto|Z-z|$ is smooth away from $z$.
Since $w-\Theta$ has a local maximum on $D$ at $Z_0$ (it has a maximum over a
neighbourhood of $Z_0$ contained in $\overline B_r(z)\subseteq D$), $\Theta$ is a
local upper test for $w$ at $Z_0$, and the subsolution inequality of
\noderef{IsInfHarmonicOn} gives $\Dinf\Theta(Z_0)\ge0$.

We compute $\Dinf\Theta(Z_0)$. Write $e=(Z_0-z)/\rho$, a unit vector. For
$g(\rho)=L\rho-\gamma\rho^2$ and $\Theta(Z)=w(z)+g(|Z-z|)$ one has
\[
  \nabla\Theta(Z_0)=g'(\rho)e=(L-2\gamma\rho)e .
\]
Since $2\gamma\rho<2\gamma r<L$, the number $L-2\gamma\rho$ is strictly
positive, so $\nabla\Theta(Z_0)$ is a positive multiple of the radial direction
$e$. The second derivative of $\Theta$ in the radial direction is
\[
  D^2\Theta(Z_0)[e]=g''(\rho)=-2\gamma .
\]
By bilinearity, $D^2\Theta(Z_0)[\nabla\Theta(Z_0)]
 =(L-2\gamma\rho)^2D^2\Theta(Z_0)[e]$, so
\[
  \Dinf\Theta(Z_0)=-2\gamma\,(L-2\gamma\rho)^2<0,
\]
contradicting $\Dinf\Theta(Z_0)\ge0$. Hence $w\le\Theta$ on
$\overline B_r(z)$, i.e.
\[
  w(Z)-w(z)\le L|Z-z|-\gamma|Z-z|^2\le L|Z-z|
       \qquad(Z\in\overline B_r(z)).
\]
Letting $\gamma\downarrow0$ is unnecessary here, since the final bound no longer
involves $\gamma$; letting $L\downarrow L_r^+(w;z)$ gives
\eqref{eq:cone-upper}.

\emph{Step 3: the lower bound.}
Apply Step 2 to $-w$, which is infinity-harmonic on $D$ by Step 1, with the same
$z$ and $r$. This gives
$(-w)(Z)-(-w)(z)\le L_r^+(-w;z)|Z-z|$ on $\overline B_r(z)$, which by
\eqref{eq:cone-neg} is exactly \eqref{eq:cone-lower}.

\emph{Step 4: monotonicity in the radius.}
Let $r'\ge r$ with $\overline B_{r'}(z)\subseteq D$; then $r'>0$. Applying Step 2
at radius $r'$ gives $w(z+h)-w(z)\le L_{r'}^+(w;z)|h|$ for all $|h|\le r'$, in
particular for all $|h|\le r$. Since $L_{r'}^+(w;z)\ge0$ by Step 0 and
$|h|\le r$,
\[
  \frac{w(z+h)-w(z)}{r}\le\frac{L_{r'}^+(w;z)|h|}{r}
       \le L_{r'}^+(w;z)\qquad(|h|\le r).
\]
Taking the supremum over $|h|\le r$ gives $L_r^+(w;z)\le L_{r'}^+(w;z)$. The
same argument applied to $-w$, together with \eqref{eq:cone-neg}, gives
$L_r^-(w;z)\le L_{r'}^-(w;z)$. This is \eqref{eq:cone-monotone}.
\end{proof}
